\section{Preliminaries}\label{section2}
In this section we shall summarize some notation and results from our papers \cite{B:S1,B:S2,B:S3}, which will be needed in the following. We use the following notation for the orthonormal polynomials
\begin{gather}%\label{eq:p}
P_n(x)=b_{n,n}x^n+b_{n-1,n}x^{n-1}+\cdots + b_{1,n}x+b_{0,n},\nonumber\\
x^n=c_{n,n}P_n(x)+c_{n-1,n}P_{n-1}(x)+\cdots + c_{1,n}P_1(x)+c_{0,n}P_0(x).\label{eq:x}
\end{gather}
By (\ref{eq:lead}) we get
\begin{gather}\label{eq:b_n} b_{n,n}={1\over b_0b_1\cdots b_{n-1}},\qquad c_{n,n}=b_0b_1\cdots b_{n-1}.
\end{gather}
The matrices $\mathcal B=\{b_{i,j}\}$ and $\mathcal C=\{c_{i,j}\}$ with the assumption
\begin{gather*}
b_{i,j}=c_{i,j}=0\qquad \text{for} \quad i>j
\end{gather*}
 are upper-triangular. Since $\mathcal B$ and $\mathcal C$ are transition matrices between two sequences of linearly independent systems of functions, we have
\begin{gather*}
\mathcal B\mathcal C=\mathcal C\mathcal B=\mathcal I.
\end{gather*}
Observe that (for say $m\le n$)
\begin{gather*}s_{m+n} = h_{m,n}=\big\langle x^m,x^n\big\rangle_{L^2(\mu)} = c_{0,m}c_{0,n}+c_{1,m}c_{1,n}+\cdots +c_{m,m}c_{m,n} =\big(\mathcal C^t \mathcal C\big)_{m,n},
\end{gather*}
so $\mathcal{H}=\mathcal C^t \mathcal C$.

Hence we have
\begin{gather}\label{eq:I=B(B^tH)}
\mathcal I=\mathcal B\mathcal C=\mathcal B\big(\mathcal B^t\mathcal C^t\mathcal C\big)=\mathcal B\big(\mathcal B^t\mathcal{H}\big).
\end{gather}
The multiplication is well defined as $\mathcal B$, $\mathcal C$ are upper-triangular matrices and $\mathcal B^t$, $\mathcal C^t$ are lower-triangular, so we are always dealing with finite sums in the matrix products.

Let us now assume that the moment problem is indeterminate. We recall from \cite[Proposition~4.2]{B:S1} that $\mathcal B$ is a Hilbert--Schmidt matrix and that \begin{gather}\label{eq:A=BBt}
\mathcal{A}=\mathcal B\mathcal B^t,
\end{gather}
 This means
\begin{gather}\label{eq:akl}
a_{k,l}=\sum^\infty_{n=\max(k,l)}b_{k,n}b_{l,n},
\end{gather}
hence
\begin{gather*}
a_{k,l}^2\le \left (\sum_{n=k}^\infty b_{k,n}^2\right ) \left (\sum_{n=l}^\infty b_{l,n}^2\right ).\end{gather*}
Let us recall some of the arguments leading to these results. Because of indeterminacy it
is known that the series
\begin{gather*}
\sum_{n=0}^\infty |P_n(z)|^2
\end{gather*}
 is convergent uniformly for $z$ in bounded subsets of $\mathbb{C}$. Moreover, the sum has subexponential growth by a result of M. Riesz, cf. \cite[p. 56]{Ak}, i.e., it is dominated by a multiple of ${\rm e}^{\varepsilon |z|}$ for any $\varepsilon>0$.
Fixing $r>0$, we have
\begin{gather*}
P_n\big(r{\rm e}^{{\rm i}t}\big)=\sum_{k=0}^n b_{k,n}r^k{\rm e}^{{\rm i}tk}.
\end{gather*}
By Parseval's identity we obtain
\begin{gather*}
{1\over 2\pi}\int_0^{2\pi}\big|P_n\big(r{\rm e}^{{\rm i}t}\big)\big|^2{\rm d}t = \sum_{k=0}^n b_{k,n}^2r^{2k}.
\end{gather*}
 Next
\begin{gather*}
{1\over 2\pi}\int_0^{2\pi}\sum_{n=0}^\infty \big|P_n\big(r{\rm e}^{{\rm i}t}\big)\big|^2{\rm d}t=\sum_{n=0}^\infty\sum_{k=0}^n b_{k,n}^2r^{2k}=\sum_{k=0}^\infty r^{2k}\sum_{n=k}^\infty b_{k,n}^2.
\end{gather*}
Therefore, denoting
\begin{gather}\label{eq:c_k}
c_k=\left (\sum_{n=k}^\infty b_{k,n}^2\right )^{1/2}=\sqrt{a_{k,k}},
\end{gather}
 we have
\begin{gather*}\label{eq:r}
\sum_{k=0}^\infty c_k^2r^{2k}={1\over 2\pi}\int_0^{2\pi}\sum_{n=0}^\infty \big|P_n\big(r{\rm e}^{{\rm i}t}\big)\big|^2{\rm d}t.
\end{gather*}
(For $r=1$ this proves that $\mathcal B$ is Hilbert--Schmidt, and in particular the series in~\eqref{eq:akl} is absolutely convergent.)

\begin{lem}\label{thm:suffAH} Consider an indeterminate moment problem. If
\begin{gather}\label{eq:suffAH}
\sum_{n,l=0}^\infty |b_{k,n}b_{l,n}s_{l+m}|<\infty\qquad \mbox{for all}\quad k,m\ge 0,
\end{gather}
then property {\rm (aci)} from Definition~{\rm \ref{thm:defAH}} holds.

For \eqref{eq:suffAH} to hold it suffices that
\begin{gather}\label{eq:suffAH1}
\sum_{l=0}^\infty c_l |s_{l+m}|<\infty\qquad\mbox{for all}\quad m\ge 0.
\end{gather}
In particular {\rm (cs*)} implies {\rm (aci)}.
\end{lem}

\begin{proof} We first notice that \eqref{eq:suffAH} implies (ac) because of \eqref{eq:akl}. We next note that \eqref{eq:suffAH} implies
\begin{gather*}
\mathcal B\big(\mathcal B^t\mathcal H\big)=\big(\mathcal B\mathcal B^t\big)\mathcal H,
\end{gather*}
so by \eqref{eq:I=B(B^tH)} and \eqref{eq:A=BBt} we get
\begin{gather*}
\mathcal I=\mathcal B\big(\mathcal B^t\mathcal H\big)=\big(\mathcal B\mathcal B^t\big)\mathcal H=\mathcal A\mathcal H.
\end{gather*}
Suppose next that \eqref{eq:suffAH1} holds. By the Cauchy--Schwarz inequality
\begin{gather*}
\sum_{n,l=0}^\infty |b_{k,n}b_{l,n}s_{l+m}|\le \sum_{l=0}^\infty c_kc_l|s_{l+m}|<\infty,
\end{gather*}
where we have used \eqref{eq:c_k}.
\end{proof}
